% Fundamental cycles from spanning tree T={12,23,34} by adding 41 or 13.
\begin{tikzpicture}[x=1mm,y=1mm,font=\figurefont\fontsize{8}{10}\selectfont,
  vertex/.style={circle,draw=FigureBlue,fill=FigureBlueFill,
    minimum size=4.6mm,inner sep=0pt,line width=.55pt},
  tree/.style={-{Stealth[length=1.65mm]},draw=FigureStroke,line width=.7pt},
  nontree/.style={-{Stealth[length=1.65mm]},draw=FigureMuted,dashed,line width=.6pt},
  cyc1/.style={-{Stealth[length=1.65mm]},draw=FigureBlue,line width=.82pt},
  cyc2/.style={-{Stealth[length=1.65mm]},draw=FigureGreen,line width=.82pt},
  value/.style={fill=white,text=FigureInk,inner sep=.7pt},
  panel/.style={font=\figurefont\fontsize{7.5}{9}\selectfont,text=FigureInk}]
  % Panel (a): spanning tree T and two non-tree edges
  \begin{scope}[xshift=0mm]
    \node[vertex] (a1) at (0,16) {$1$};
    \node[vertex] (a2) at (18,16) {$2$};
    \node[vertex] (a3) at (18,0) {$3$};
    \node[vertex] (a4) at (0,0) {$4$};
    \draw[tree] (a1) -- (a2);
    \draw[tree] (a2) -- (a3);
    \draw[tree] (a3) -- (a4);
    \draw[nontree] (a4) -- (a1);
    \draw[nontree] (a1) -- (a3);
    \node[panel] at (9,-11) {(a) 生成树 $T$};
  \end{scope}

  % Panel (b): add non-tree edge 41 -> c_1 = (1,1,1,1,0)^T
  \begin{scope}[xshift=36mm]
    \node[vertex] (b1) at (0,16) {$1$};
    \node[vertex] (b2) at (18,16) {$2$};
    \node[vertex] (b3) at (18,0) {$3$};
    \node[vertex] (b4) at (0,0) {$4$};
    \draw[cyc1] (b1) -- node[value,above] {$+1$} (b2);
    \draw[cyc1] (b2) -- node[value,right] {$+1$} (b3);
    \draw[cyc1] (b3) -- node[value,below] {$+1$} (b4);
    \draw[cyc1] (b4) -- node[value,left] {$+1$} (b1);
    \draw[nontree] (b1) -- node[value,above right] {$0$} (b3);
    \node[panel] at (9,-11) {(b) 添 $41$ 得 $\bm c_1$};
  \end{scope}

  % Panel (c): add non-tree edge 13 -> c_2 = (-1,-1,0,0,1)^T
  \begin{scope}[xshift=72mm]
    \node[vertex] (c1) at (0,16) {$1$};
    \node[vertex] (c2) at (18,16) {$2$};
    \node[vertex] (c3) at (18,0) {$3$};
    \node[vertex] (c4) at (0,0) {$4$};
    \draw[cyc2] (c1) -- node[value,above] {$-1$} (c2);
    \draw[cyc2] (c2) -- node[value,right] {$-1$} (c3);
    \draw[nontree] (c3) -- node[value,below] {$0$} (c4);
    \draw[nontree] (c4) -- node[value,left] {$0$} (c1);
    \draw[cyc2] (c1) -- node[value,above right] {$+1$} (c3);
    \node[panel] at (9,-11) {(c) 添 $13$ 得 $\bm c_2$};
  \end{scope}
\end{tikzpicture}
